\section{Memory Management and File System} 
\label{sec:mem}

A central part of Apollo-RTOS is its memory allocation and management system. All
runtime allocation is handled through instances of a pool allocator, which provides a
simple and efficient form of dynamic memory management.
\begin{emptythm}
    \begin{itemize}
        \item The allocator maintains a linked list of free memory chunks.
        \item On an allocation request, it searches the free list for a chunk of
            sufficient size.
        \item If a suitable chunk is found, it is returned; otherwise, more memory is
            requested from the global allocator.
        \item When \lstinline[style=myagc]|free| is called, the chunk is returned to
            the free list.
    \end{itemize}
\end{emptythm}

\begin{lstlisting}[
    style=myagc,
    title={\file{include/utility/allocator.h}},
    linebackgroundcolor={\line{12} \line{16} \line{19} \line{22}}
]
template <byte_t *(*alloc_func)(size_t), size_t alignment>
class allocator
{
  chunk_t free_hd; 
  static constexpr size_t header_sz = sizeof(chunk_t);

public:
  // allocate chunk of sz bytes
  byte_t *alloc(size_t sz)
  {
    sz = roundup(sz, alignment);
    auto chnk_sz = header_sz + sz;

    chunk_t *chunk = /* find a chunk in free_hd, nullptr otherwise */;
    if (chunk == nullptr) {
      chunk = (chunk_t *)alloc_func(chnk_sz); // get new chunk
      chunk->sz = sz;
    } else {
      /* remove chunk from free_hd */
    }

    return (byte_t *)chunk + header_sz;
  }

  // return chunk allocated by this allocator
  void dealloc(byte_t *ptr)
  {
    if (ptr == nullptr)
      return;
    auto chunk = (chunk_t *)(ptr - header_sz);
    /* insert chunk to free_hd */
  }
};
\end{lstlisting}

Having an instantiable allocator instead of a global \lstinline[style=myagc]|malloc|
function has key benefits:
\begin{itemize}
    \item Since the hardware lacks a Memory Protection Unit (MPU), this helps isolate
        unrelated memory allocations in software, reducing memory corruption risks.
    \item It provides a convenient interface for creating local pool allocators.
\end{itemize}

Using this allocator, we implement kernel-space memory management via
\lstinline[style=myagc]|kmalloc| and \lstinline[style=myagc]|kfree|.
\begin{lstlisting}[
    style=myagc, 
    title={\file{src/core/memory.cpp}}
]
namespace kmem
{
allocator<alloc_heap, 8> kpool;

void *kmalloc(size_t sz) { return kpool.alloc(sz); }

void kfree(void *ptr)
{
  if (ptr == nullptr)
    return;
  kpool.dealloc((byte_t *)ptr);
}
} // namespace kmem
\end{lstlisting}

For user-space memory allocation, the implementation is more involved. When a process
allocates memory, it is expected to free it explicitly. However, if the process
encounters an error or is terminated prematurely, the scheduler must clean up any
unfreed memory.

\begin{emptythm}
    Each process descriptor maintains a list of allocated memory chunks. When a process
    terminates, all remaining allocations in this list are freed automatically.
\end{emptythm}

\begin{lstlisting}[
    style=myagc, 
    title={\file{src/core/memory.cpp}}
]
namespace heap
{
allocator<alloc_heap, 8> pool;

void *malloc(size_t sz)
{
  byte_t *ptr = pool.alloc(sz);
  /* insert chunk to the current process' allocated list */
  return ptr;
}

void free(void *ptr)
{
  if (ptr == nullptr)
    return;
  /* remove chunk from process' allocated list */
  pool.dealloc((byte_t *)ptr);
}

// called at process teardown
void cleanup(chunk_t *hd)
{
  while (hd->next != nullptr) {
    auto ptr = (byte_t *)hd->next + sizeof(chunk_t);
    /* remove hd->next from the list */
    pool.dealloc(ptr);
  }
}
} // namespace heap
\end{lstlisting}


\subsection{File System}

\paragraph*{Flash} The \texttt{micro:bit} features approximately 256KB of flash memory.
A build of Apollo-RTOS with the \texttt{-O2} flag produces a binary of about 92KB,
while using \texttt{-Os} reduces it to 76KB. In either case, there is ample remaining
space in flash that can be used for long-term storage. However, flash memory has
significant drawbacks:
\begin{itemize}
    \item Inefficient writes: Writing to flash requires first erasing entire 1KB pages,
        taking about \textbf{22.3ms} per page. Writing individual bytes afterward takes
        \textbf{20$\mu$s} each \citep{board_spec}.
    \item Limited endurance: The \texttt{micro:bit}'s flash memory has an endurance of
        only \textbf{20,000 erase cycles} and retains data for approximately \textbf{10
        years} under normal conditions. Excessive writes degrade reliability.
\end{itemize}
Despite these limitations, reading from flash is fast since the CPU has direct access
to it—after all, the program itself runs from flash. This makes flash memory suitable
for read-heavy, infrequent-write storage needs.

\paragraph*{FRAM} Ferroelectric RAM (FRAM) addresses flash memory's shortcomings.
Unlike flash, FRAM is non-volatile, with significantly higher endurance and retention
rates. The MB85RC2567 FRAM chip \citep{fram} used in Apollo-RTOS provides:
\begin{itemize}
    \item \textbf{\(10^{12}\) erase cycles} -- far exceeding flash longevity.
    \item \textbf{95+ years} of data retention under normal conditions.
    \item 32KB capacity, operating at 1MHz I2C, though we use 100KHz.
\end{itemize}
However, FRAM has its own trade-offs:
\begin{itemize}
    \item I2C communication overhead: Since the \texttt{micro:bit} lacks integrated
        FRAM, it is connected as an external I2C peripheral. This introduces bus
        contention, affecting read/write performance.
    \item Slower reads than flash: Unlike flash, which the CPU can access directly,
        FRAM requires I2C transactions, making reads slower. However, writes are
        significantly faster than flash.
\end{itemize}

Since both storage options have their strengths and weaknesses, Apollo-RTOS implements
file systems on both. Instead of maintaining two separate file systems, we designed an
abstract block device that can be instantiated with different parameters for each
storage medium.

A block device is characterized by:
\begin{itemize}
    \item Constants defining block size, total blocks, etc.
    \item Read/Write methods for interacting with the device.
\end{itemize}
Since these constants are frequently used in file system operations, we implemented
them as template parameters rather than function parameters. This allows the compiler
to optimize operations more effectively.

Thus, our file system implementation is entirely header-only, residing in the
\texttt{include/fs/} directory.


\subsection{File system implementation}

The file system consists of two layers of tables with extensive bookkeeping. Below, we
outline key concepts and methods.

\paragraph*{\texttt{fs\_desc\_t}} defines the file system parameters. Two instances—one
for flash, another for FRAM—are constructed at compile time and passed as template
parameters to the implementation.

\begin{lstlisting}[
    style=myagc,
    title={\file{include/fs/fs\_desc.h}},
    linebackgroundcolor={\lines{18}{19} \lines{4}{5} \line{8} \line{10}}
]
template <...>
struct fs_desc_t {
  using addr_t = /* type of address supported by device */;
    // for flash, it is uint32_t since writes must be word aligned
    // for fram, it can be uint16_t to save space 

  size_t blk_sz; /* size of a block */
    // for flash, it is 1024, 
    //   since that's the smallest unit we can effectively write
    // for FRAM, it is 128

  size_t max_num_blks; /* man number of blocks per file */
  size_t n_inodes; /* number of files */

  size_t n_open_files; /* number of open files */

  // methods to transfer data to and from the target device
  void (*write)(addr_t addr, uint8_t *buf, size_t buf_sz);
  void (*read)(addr_t addr, uint8_t *buf, size_t buf_sz);
};
\end{lstlisting}


\subsubsection{File system backend}

Each file is represented by an \lstinline[style=myagc]|inode_t|, inspired by Linux
inodes, containing file metadata.
\begin{lstlisting}[
    style=myagc,
    title={\file{include/fs/fs\_bck.h}},
    linebackgroundcolor={}
]
template <typename desc_t, desc_t desc>
struct inode_t {
  using addr_t = typename desc_t::addr_t;
  using blk_addr_t = uint8_t; 
    // in either flash and fram, we can refer to a block by a 8 bit uint

  size_t nblks; /* number of blocks in this file */
  blk_addr_t blks[desc.max_num_blks]; /* blocks in this file */

  uint32_t flag; /* stats encoded into a single uint32_t */
};
\end{lstlisting}

We use \lstinline[style=myagc]|fn_t = uint8_t| as the file name. We currently don't map
file name strings to the corresopnding inode, and directly refer the file by the inode
number.

The file system header maintains:
\begin{itemize}
    \item A table of inodes.
    \item A \lstinline[style=myagc]|bitset| tracking free and occupied blocks.
\end{itemize}

\begin{lstlisting}[
    style=myagc,
    title={\file{include/fs/fs\_bck.h}},
    linebackgroundcolor={}
]
template <typename desc_t, desc_t desc>
struct fs_hd_t {
  inode_t inode_tbl[desc.n_inodes];
  bitset<N_BLKS> free_set; /* N_BLKS is total number of available blks */

  bool valid() const;
  void format(); // format the file system

  // find the inode for the file fn
  inode_t &find(fn_t fn);

  // find nblks unclaimed blks and assign them to the given buffer
  void alloc_blks(blk_addr_t *buf, size_t nblks);

  // reclaim the given blks
  void dealloc_blks(blk_addr_t *buf, size_t nblks);
};
\end{lstlisting}


The file system backend is then a wrapper around the header, that provides methods for:
\begin{itemize}
    \item \lstinline[style=myagc]|open| and \lstinline[style=myagc]|remove| files.
    \item \lstinline[style=myagc]|load| and \lstinline[style=myagc]|store| data.
    \item \lstinline[style=myagc]|mount| and \lstinline[style=myagc]|format|
        operations.
\end{itemize}

\begin{lstlisting}[
    style=myagc,
    title={\file{include/fs/fs\_bck.h}},
    linebackgroundcolor={}
]
template <typename desc_t, desc_t desc>
class fs_t
{
  fs_hd_t fs_hd;

public:
  void format();
  void mount(); /* if the header is not valid, it also formats it */
  void umount() const; 

  // Open a file using Linux like flags O_CREATE, O_WRITE, O_SHARED etc.
  const inode_t *open(fn_t fn, size_t sz, uint32_t flags);

  void remove(fn_t fn);

  // Load the contents of the file to the given buffer
  void load(const inode_t *inode, uint8_t *buf, size_t buf_sz) const;

  // Store the contents of the given buffer to the file
  void store(const inode_t *inode, uint8_t *buf, size_t buf_sz) const;

  // Store the contents of the buffer at a particular offset.
  void store(const inode_t *inode, size_t off, uint8_t *buf,
             size_t buf_sz) const;
};
\end{lstlisting}

\lstinline[style=myagc]|load| and \lstinline[style=myagc]|store| map buffer contents to
blocks, allowing files to have non-contiguous block allocations.

% FIXME TODO in the future maybe?
Currently a file's size cannot be modified. But it should be fairly simple to increase
the size of a file, just by allocating more blocks to it. 


\subsubsection{File System Frontend}

The frontend maintains a table of open files and maps them to runtime memory addresses.
This mapping can be either private or shared:
\begin{itemize}
    \item Private mapping: A file is mapped to a runtime buffer local to the process.
    \item Shared mapping: Multiple processes access the same runtime representation of
        the file, enabling inter-process communication (IPC).
\end{itemize}
This mechanism is inspired by the \lstinline[style=myagc]|mmap| function in Linux.

File operations occur in the following stages:

\paragraph*{Opening a File} When a process requests to open a file \texttt{fn}:
\begin{itemize}
    \item The frontend retrieves the associated inode, creating one if it doesn't
        exist.
    \item It locates an available file descriptor in the open files table.
    \item A file handler is wrapped around the descriptor and returned to the process.
\end{itemize}

\paragraph*{Mapping a File} The mapping strategy depends on whether the
\lstinline[style=myagc]|O_SHARED| flag is set during the \lstinline[style=myagc]|open|
call:
\begin{itemize}
    \item Shared mapping:
        \begin{itemize}
            \item A buffer is allocated, and its address is stored in the file
                descriptor.
            \item If a mapping already exists, it is reused.
            \item All \lstinline[style=myagc]|load| and \lstinline[style=myagc]|store|
                operations use this shared buffer.
        \end{itemize}
    \item Private mapping:
        \begin{itemize}
            \item A buffer is allocated and stored in the file handler.
            \item Other processes cannot access this mapping.
        \end{itemize}
\end{itemize}

\paragraph*{Loading and Storing a File} Once mapping is established, all
\lstinline[style=myagc]|load| and \lstinline[style=myagc]|store| operations reference
the appropriate buffer from either the file descriptor (shared mapping) or file handler
(private mapping).

The file descriptor maintains:
\begin{itemize}
    \item A pointer to the underlying inode.
    \item A reference count for read and write operations.
    \item Information on shared memory mappings.
\end{itemize}

\begin{lstlisting}[
    style=myagc,
    title={\file{include/fs/fs\_impl.h}},
    linebackgroundcolor={}
]
class fd_t
{
  uint8_t w_cnt, r_cnt; /* read and write reference counter */

public:
  fn_t fn;
  mutable mutex<8> mtx = {"fd"}; /* file operations should be atomic */
  const inode_t *inode; /* the underlying file */
  mmap_unit_t mu; /* mapping information */

public:
  // if this file descriptor refers the file fn, increse ref count
  // otherwise mark this descriptor in use
  void acquire(fn_t _fn, const inode_t *_inode, bool w_en);

  // decrease ref count, 
  // if it was the last reference to the file, close this descriptor
  void release(bool w_en);
};
\end{lstlisting}

The file handler is what the process interacts with. It holds a pointer to the file
descriptor and provides methods for interacting with the file, ensuring smooth access
and management.

\begin{lstlisting}[
    style=myagc,
    title={\file{include/fs/fs\_impl.h}},
    linebackgroundcolor={}
]
class file_t
{
  fd_t *fd;
  mmap_unit_t mu; /* private mapping information */
  const bool is_shared; /* is the mapping shared? */
  const bool w_en = false; /* have write permission? */

public:
  // open the file in an RAII style, close it in the desctructor
  file_t(fn_t fn, size_t _sz, uint32_t flags);

  ~file_t();
  void close();

  // create runtime mapping
  void *mmap(size_t sz);

  // instead of allocating new memory, use the given memory address 
  void mmap(uint8_t *buf, size_t sz);

  // remove the mapping
  void unmap();

  void load();
  void store();
  // store a chunk of size sz at offset off
  void store(size_t off, size_t sz);
};
\end{lstlisting}

Finally, the top-level file system implementation maintains:
\begin{itemize}
    \item An array of open files.
    \item An \lstinline[style=myagc]|fs_t| instance referring to the backend.
    \item Methods for interacting with the file system backend.
\end{itemize}

\begin{lstlisting}[
    style=myagc,
    title={\file{include/fs/fs\_impl.h}},
    linebackgroundcolor={}
]
template <desc_t desc>
class fs_impl_t
{
  // the backend implementation
  inline static fs_t fs;

  // pool allocator providing memory for shared mapping
  inline static allocator<alloc_heap, 4> pool; 

  // table of open files
  inline static fd_t open_files[desc.n_open_files] = {};

public:
  static file_t open(fn_t fn, size_t sz, uint32_t flags);
  static void remove(fn_t fn);
  static void mount();
  static void format();
};
\end{lstlisting}


\subsubsection{File system instantiation}

So far, the file system has been an abstract concept implemented through a complex
network of templates. To create actual file systems, we need to instantiate these
templates with specific parameters. After defining the required parameters for both the
flash and FRAM file systems, the instantiation process becomes straightforward:
\begin{lstlisting}[
    style=myagc,
    title={\file{include/fs/fs.h}},
    linebackgroundcolor={}
]
// instantiate with the parameters for the flash fs
class flash : public fs_impl_t<_flash::desc> { };

// the same for the fram fs
class fram : public fs_impl_t<_fram::desc> { };
\end{lstlisting}


Now, we can use files as follows:

\begin{lstlisting}[
    style=myagc,
    title={\file{apollo/test.cpp}},
]
void file_rw_test()
{
  fn_t fn = 10;

  auto file = fram::open(fn, 16, O_WRITE | O_CREATE);
  auto t = (char *)file.mmap(16);

  const char *s = "It IS a string";
  while (*s != '\0')
    *t++ = *s++;
  *t = '\0';

  file.store();
  // the file is automatically closed by the destructor
}
\end{lstlisting}

